% ======================================================================
\section{2.2 Tożsamość Schönhagego}
% ======================================================================

\begin{frame}{Dwa iloczyny na rozłącznych zmiennych}
  \begin{columns}[T]
    \begin{column}{0.5\textwidth}
      \textbf{\Zew{Iloczyn zewnętrzny}} wektorów $(x_1,x_2,x_3)$ i $(y_1,y_2,y_3)$:
      \[\text{dziewięć liczb } x_iy_j,\quad 1\le i,j\le3.\]
      Jest to iloczyn macierzy $3\times1$ przez $1\times3$.
      \begin{center}
        \begin{tikzpicture}[scale=0.32]
          \begin{scope}[shift={(0,0)}]\draw[fill=tlozew] (0,0) rectangle (1,3); \draw (0,0) grid (1,3);\end{scope} \node at (0.5,-0.5) {\scriptsize $3{\times}1$};
          \node at (1.6,1.5) {$\cdot$};
          \begin{scope}[shift={(2.2,1)}]\draw[fill=tlozew] (0,0) rectangle (3,1); \draw (0,0) grid (3,1);\end{scope} \node at (3.7,0.5) {\scriptsize $1{\times}3$};
          \node at (5.8,1.5) {$=$};
          \begin{scope}[shift={(6.4,0)}]\draw[fill=tlozew] (0,0) rectangle (3,3); \draw (0,0) grid (3,3);\end{scope} \node at (7.9,-0.5) {\scriptsize $3{\times}3$};
        \end{tikzpicture}
      \end{center}
    \end{column}
    \begin{column}{0.5\textwidth}
      \textbf{\Wew{Iloczyn skalarny}} wektorów $(p_{11},p_{12},p_{21},p_{22})$ i $(q_{11},q_{12},q_{21},q_{22})$:
      \[\textstyle\sum_{i,j=1}^{2}p_{ij}q_{ij}.\]
      Jest to iloczyn macierzy $1\times4$ przez $4\times1$.
      \begin{center}
        \begin{tikzpicture}[scale=0.32]
          \begin{scope}[shift={(0,1)}]\draw[fill=tlowew] (0,0) rectangle (4,1); \draw (0,0) grid (4,1);\end{scope} \node at (2,0.5) {\scriptsize $1{\times}4$};
          \node at (4.6,1.5) {$\cdot$};
          \begin{scope}[shift={(5.2,-0.5)}]\draw[fill=tlowew] (0,0) rectangle (1,4); \draw (0,0) grid (1,4);\end{scope} \node at (5.7,-1) {\scriptsize $4{\times}1$};
          \node at (6.8,1.5) {$=$};
          \draw[fill=tlowew] (7.4,1) rectangle (8.4,2); \node at (7.9,0.5) {\scriptsize $1{\times}1$};
        \end{tikzpicture}
      \end{center}
      {\scriptsize Indeksy $\{1,2\}^2$ zamiast $\{1,2,3,4\}$ ułatwią zapis $\ph$, $\qh$.}
    \end{column}
  \end{columns}
  \pause
  \begin{itemize}\small
    \item Osobno: $9+4=\mathbf{13}$ mnożeń.
    \item<3-> \textbf{Schönhage [Sch81]:} wystarczy \alert{\textbf{10}} mnożeń, jeśli dopuścimy pewne
      wyrazy błędu (okażą się nieszkodliwe).
  \end{itemize}
\end{frame}

\begin{frame}{Zapis trójliniowy: wielomian $G$}
  \begin{itemize}
    \item \textbf{7 zmiennych lewych:} $\Zew{x_1,x_2,x_3}$, $\Wew{p_{11},p_{12},p_{21},p_{22}}$.
    \item \textbf{7 zmiennych prawych:} $\Zew{y_1,y_2,y_3}$, $\Wew{q_{11},q_{12},q_{21},q_{22}}$.
    \item \textbf{10 zmiennych wyjściowych:} $\Zew{z_{ij}}$ ($1\le i,j\le3$) i $\Wew{z_0}$, po jednej na każdą z dziesięciu liczb do policzenia.
  \end{itemize}
  \pause
  \[
    G:=\sum_{i,j=1}^{3}\Zew{x_iy_j\,z_{ij}}\;+\;\Big(\sum_{i,j=1}^{2}\Wew{p_{ij}q_{ij}}\Big)\Wew{z_0}.
  \]
  \pause
  \begin{itemize}
    \item Współczynnik przy $z_{ij}$ to $x_iy_j$ (iloczyn zewnętrzny).
    \item Współczynnik przy $z_0$ to $\sum p_{ij}q_{ij}$ (iloczyn skalarny).
    \item Chcemy zapisać $G$ (z dokładnością do błędu) jako sumę \textbf{dziesięciu} wyrazów
      postaci (forma liniowa w zmiennych lewych)$\times$(forma w prawych)$\times$(forma w wyjściowych).
  \end{itemize}
\end{frame}

\begin{frame}{Macierze $\ph$ i $\qh$}
  Rozszerzamy macierze $2\times2$ zmiennych $p$ i $q$ do $3\times3$ tak, żeby
  \textbf{każda kolumna $\ph$} i \textbf{każdy wiersz $\qh$} sumowały się do zera:
  \[
    \ph=\begin{pmatrix} p_{11}&p_{12}&0\\ p_{21}&p_{22}&0\\ -p_{11}-p_{21}&-p_{12}-p_{22}&0\end{pmatrix},
    \qquad
    \qh=\begin{pmatrix} q_{11}&q_{12}&-q_{11}-q_{12}\\ q_{21}&q_{22}&-q_{21}-q_{22}\\ 0&0&0\end{pmatrix}.
  \]
  \pause
  \begin{itemize}
    \item Elementy $\ph_{ij}$, $\qh_{ij}$ to \emph{formy liniowe}, a nie nowe zmienne.
    \item $\displaystyle\sum_{i,j=1}^{3}\ph_{ij}\qh_{ij}=\sum_{i,j=1}^{2}p_{ij}q_{ij}$, bo trzecia
      kolumna $\ph$ i trzeci wiersz $\qh$ są zerowe, a dla $i,j\le2$ mamy $\ph_{ij}=p_{ij}$, $\qh_{ij}=q_{ij}$.
  \end{itemize}
\end{frame}

\begin{frame}{Dziesięć składników i ich formy liniowe}
  Składniki to $P_{ij}$ (dla $1\le i,j\le3$, odpowiadające wpisom $\ph$ i $\qh$) oraz $P_0$.
  Każdy \textbf{składnik} $\lambda$ jest iloczynem trzech form liniowych:
  $\varphi_\lambda$ w zmiennych lewych, $\psi_\lambda$ w prawych, $\chi_\lambda$ w wyjściowych.
  \vspace{0.4em}
  \begin{center}
    \renewcommand{\arraystretch}{1.4}
    \begin{tabular}{c|c|c|c}
      składnik $\lambda$ & $\varphi_\lambda$ (lewe) & $\psi_\lambda$ (prawe) & $\chi_\lambda$ (wyjściowe)\\\hline
      $P_{ij}$, $1\le i,j\le3$ & $x_i+\ph_{ij}$ & $y_j+\qh_{ij}$ & $z_{ij}+z_0$\\
      $P_0$ & $-(x_1+x_2+x_3)$ & $y_1+y_2+y_3$ & $z_0$
    \end{tabular}
  \end{center}
  \pause
  \begin{itemize}
    \item 10 składników to \textbf{10 mnożeń} $\varphi_\lambda\cdot\psi_\lambda$.
    \item Każde $\varphi_\lambda$ i każde $\psi_\lambda$ ma \textbf{co najwyżej 3} jednomiany o współczynnikach $\pm1$.
  \end{itemize}
\end{frame}

\begin{frame}{Lemat 6: tożsamość Schönhagego}
  \begin{lemat}{6}
    Sumując po dziesięciu składnikach,
    \[
      \sum_\lambda\varphi_\lambda\psi_\lambda\chi_\lambda=G+E,
      \qquad\text{gdzie}\qquad
      E:=\sum_{i,j=1}^{3}\big(x_i\qh_{ij}+\ph_{ij}y_j+\ph_{ij}\qh_{ij}\big)z_{ij}.
    \]
  \end{lemat}
  \pause
  \begin{itemize}
    \item Lewa strona: 10 wyrazów, czyli \textbf{10 mnożeń} $\varphi_\lambda\cdot\psi_\lambda$.
    \item Współczynnik przy $z_0$ po lewej to \textbf{dokładnie} iloczyn skalarny $\sum p_{ij}q_{ij}$.
    \item Współczynnik przy $z_{ij}$ to $x_iy_j$ \textbf{plus błąd} $x_i\qh_{ij}+\ph_{ij}y_j+\ph_{ij}\qh_{ij}$.
    \item<3-> Lemat 6 to kluczowa tożsamość: liczy chciane $G$ plus wielomian błędu $E$.
  \end{itemize}
\end{frame}

\begin{frame}{Dowód Lematu 6}
  \begin{proof}
    Porównujemy współczynniki przy każdej zmiennej wyjściowej po lewej stronie.
    \begin{itemize}
      \item<2-> \textbf{Przy $z_{ij}$:} $z_{ij}$ występuje tylko w $\chi_{P_{ij}}$, więc współczynnik wynosi
        \[\varphi_{P_{ij}}\psi_{P_{ij}}=(x_i+\ph_{ij})(y_j+\qh_{ij})
          =x_iy_j+\big(x_i\qh_{ij}+\ph_{ij}y_j+\ph_{ij}\qh_{ij}\big).\]
      \item<3-> \textbf{Przy $z_0$:} $z_0$ występuje we wszystkich $\chi_\lambda$, więc współczynnik wynosi
        \[
          \sum_{i,j=1}^{3}(x_i+\ph_{ij})(y_j+\qh_{ij})-(x_1+x_2+x_3)(y_1+y_2+y_3)
          =\underbrace{\sum_{i}x_i\sum_{j}\qh_{ij}}_{=0\ (\text{wiersze }\qh)}
           +\underbrace{\sum_{j}y_j\sum_{i}\ph_{ij}}_{=0\ (\text{kolumny }\ph)}
           +\sum_{i,j=1}^{3}\ph_{ij}\qh_{ij}.
        \]
        Iloczyny $x_iy_j$ się skracają, a $\sum_{i,j=1}^{3}\ph_{ij}\qh_{ij}=\sum_{i,j=1}^{2}p_{ij}q_{ij}$. \qedhere
    \end{itemize}
  \end{proof}
\end{frame}

\begin{frame}{Obserwacja 1: zmienne wewnętrzne i zewnętrzne}
  \begin{definicja}
    Zmienne iloczynu skalarnego, czyli $\Wew{p_{ij},q_{ij},z_0}$, nazywamy \Wew{\textbf{wewnętrznymi}}.
    Zmienne iloczynu zewnętrznego, czyli $\Zew{x_i,y_j,z_{ij}}$, nazywamy \Zew{\textbf{zewnętrznymi}}.
  \end{definicja}
  \pause
  \begin{obserwacja}
    \begin{itemize}
      \item Każdy jednomian $G$ jest \emph{czysto} wewnętrzny ($\Wew{p_{ij}q_{ij}z_0}$) albo
        \emph{czysto} zewnętrzny ($\Zew{x_iy_jz_{ij}}$).
      \item Każdy jednomian $E$ ma \Zew{zewnętrzną zmienną wyjściową $z_{ij}$} razem
        z \Wew{wewnętrzną zmienną lewą lub prawą}:
        $\Zew{x_i}\Wew{\qh_{ij}}\Zew{z_{ij}}$, \ $\Wew{\ph_{ij}}\Zew{y_jz_{ij}}$, \ $\Wew{\ph_{ij}\qh_{ij}}\Zew{z_{ij}}$.
    \end{itemize}
  \end{obserwacja}
  To pozwoli pokazać, że $E$ jest nieszkodliwe (Lemat 9).
\end{frame}

\begin{frame}{Obserwacja 2: kto wnosi wkład do którego wyjścia}
  \begin{columns}[T]
    \begin{column}{0.5\textwidth}
      \small
      \begin{definicja}
        Składnik $\lambda$ \textbf{wnosi wkład} do zmiennej wyjściowej $z$, jeśli $z$ występuje w $\chi_\lambda$.
      \end{definicja}
      \begin{itemize}
        \item<2-> \textbf{Każdy} składnik wnosi wkład do $z_0$.
        \item<2-> Do $z_{ij}$ wnosi wkład \textbf{tylko} $P_{ij}$.
        \item<3-> Formy $\chi_\lambda$ są sumami zmiennych bez współczynników, więc współczynnik przy $z$ po lewej stronie to
          \[\sum_{\lambda\text{ wnosi wkład do }z}\varphi_\lambda\psi_\lambda.\]
        \item<4-> To, że do $z_{ij}$ wnosi wkład tylko jeden składnik, pozwoli
          \alert{pominąć wiele mnożeń}, gdy chcemy tylko niektórych wyjść.
      \end{itemize}
    \end{column}
    \begin{column}{0.48\textwidth}
      \centering
      \begin{tikzpicture}[
          t/.style={draw,rounded corners=1pt,fill=red!10,font=\tiny,minimum width=0.5cm,inner sep=1.5pt},
          o/.style={draw,rounded corners=1pt,fill=green!12,font=\tiny,minimum width=0.5cm,inner sep=1.5pt}]
        \foreach \n [count=\k] in {11,12,13,21,22,23,31,32,33} {
          \node[t] (P\n) at ({(\k-1)*0.62},2) {$P_{\n}$};
          \node[o] (z\n) at ({(\k-1)*0.62},0) {$z_{\n}$};
          \draw[thick] (P\n) -- (z\n);
        }
        \node[t] (P0) at (5.9,2) {$P_0$};
        \node[o] (z0) at (5.9,0) {$z_0$};
        \foreach \n in {11,12,13,21,22,23,31,32,33} \draw[szary] (P\n.south) -- (z0.north);
        \draw[thick] (P0) -- (z0);
        \node[font=\scriptsize] at (2.9,2.6) {składniki $\lambda$};
        \node[font=\scriptsize] at (2.9,-0.6) {zmienne wyjściowe};
      \end{tikzpicture}
    \end{column}
  \end{columns}
\end{frame}
